
% ===================================== [APPENDIX] ====================================

\section{The leakage certificate: hypotheses, protocol, and the anatomy of the threshold}
\label{app:hyp}

Theorem~\ref{thm:leakage} holds under three hypotheses. Each is a statement about the pipeline that
produced the reconstruction rather than about the theorem, so each has to be checked against the
pipeline. We check all
three here, with numbers, for the sweep of Sec.~\ref{sec:frontier}, and add the two protocol
declarations and the decomposition that section quotes.

\paragraph{(H1) Common origin.} $A$ and $A_P$ must be referred to the same $E_0$. A shift common to
both cancels exactly---verified to $1.000000$ for shifts of $10^{-3}$, $10^{-2}$ and
$10^{-1}\,t$---which is how the hypothesis is satisfied here, a single $E_0$ being used for both
objects. If $A_P$ is instead referred to $E_0+\delta$, the discrepancy acquires an additive
$\lVert\phi_P\rVert^{2}(4/\pi)\arctan(\lvert\delta\rvert/2\eta)$, verified to $2\times10^{-5}$ in all
$48$ cases tested, \emph{which the certificate does not control}. A ground state obtained by
Rayleigh--Ritz on $60\%$ of the $N$-electron sector at $L=8$ gives $\delta=6.7\times10^{-4}\,t$, which
adds $0.0024$ to a bound of $1.2036$---$+0.20\%$---while multiplying the \emph{measured} error by
$16$; and the dependence is not even monotone, a shift of $10^{-2}\,t$ lowering the measured error at
$L=6$, fraction $0.85$, $\eta=0.18\,t$ from $2.7\times10^{-2}$ to $1.2\times10^{-2}$. The measured
errors quoted beside a bound in Sec.~\ref{sec:frontier} are therefore those of a pipeline handed the
exact $E_0$, and they bound the error of a pipeline that estimates it from neither side. No cheap a
posteriori control on $\delta$ is available: residual (Weinstein) and Temple
bounds~\cite{weinstein1934,temple1928} computed on
ground-state subspaces of $20$--$80\%$ of the $N$ sector at $L=6$ give $\arctan$ terms of $1.76$,
$1.64$, $1.40$ and $0.92$, all vacuous against the ceiling of $2$. In a shot-based pipeline $E_0$ is
itself estimated on whatever subspace the shots populated, and the certificate has never been tested
that way.

\paragraph{(H2) Normalization and support of the $L_1$ norm.} The bound is on
$\int_{\mathbb R}\lvert A-A_P\rvert\,d\omega$ normalized by
$\lVert\phi\rVert^{2}=\int_{\mathbb R}A\,d\omega$. A relative error normalized instead by $\int_W A$
over a finite window $W$ and evaluated by finite quadrature---the convention of the repository, in
which Sec.~\ref{sec:prices} reports---is a different quantity, differing by the product of
$\int_W\lvert A-A_P\rvert/\int_{\mathbb R}\lvert A-A_P\rvert$ and $\lVert\phi\rVert^{2}/\int_W A$.
Both factors are measured here: the second is $1.0100$, $1.0309$ and $1.0373$ at $\eta=0.05$, $0.15$
and $0.18\,t$ at $L=6$ and $1.0110$, $1.0278$ and $1.0335$ at $L=8$, while the first lies between
$1.0040$ and $1.0115$ across the $204$ configurations of this study. Every number in
Sec.~\ref{sec:frontier} is in the normalization of Theorem~\ref{thm:leakage}, and the two conventions
are nowhere interchanged.

\paragraph{(H3) Exact reference.} $A$ denotes the exact spectral function of the full $(N\pm1)$
sector. If the reference is itself a Lanczos reconstruction of finite depth $n_\ell$, the theorem does
\emph{not} bound $\lVert A^{(n_\ell)}-A_P^{(n_\ell)}\rVert_1$, and a measured error below the
reference's own drift measures the reference rather than the reconstruction. That drift runs from
$2.6\times10^{-5}$ ($L=12$, $\eta=0.10\,t$) to $2.8\times10^{-7}$ ($\eta=0.18\,t$), and at $L=10$ the
reorthogonalization gap, $1.71\times10^{-5}$, is of the same order as the depth drift. We therefore
quote no measured error anywhere in Sec.~\ref{sec:frontier} within a factor of $20$ of its own
reference drift: twelve rows at $L\ge10$ fall under that rule and are left blank rather than rounded,
the weakest surviving cell---$L=12$, $\eta=0.10\,t$---standing at a factor $47$. One gap is declared
and not bounded. The reconstructions this paper publishes are produced by bare Lanczos at depth $150$,
while the certificate is evaluated on the reorthogonalized Krylov space; the difference between the
two has been measured at a single point at $L=8$---relative $L_1$ of $1.2454\times10^{-3}$ against
$1.2464\times10^{-3}$---and never bounded, and at $L=12$ it has never been measured at all.

\paragraph{Ranking, ties and support.} Two properties of the cut are declared with the protocol. Only the reflection-forbidden determinants ($4$, $18$, $48$, $200$ at $L=6$--$12$) carry exactly zero accumulated Born weight, and they rank last, so the cut is never filled with configurations the
device would never sample in the infinite-shot limit; ties at the cut involve two determinants (four
at $L=6$, at fractions $0.80$ and $0.85$) and move the bound by at most $6.6\times10^{-4}$ relative.
Above the support floor $\lvert\operatorname{supp}(\phi)\rvert/D$ of Sec.~\ref{sec:retraction}, on the other hand, the
ranking selects among determinants outside $\operatorname{supp}(\phi)$: at fraction $0.85$ that is
$30.2\%$ of $S$ at $L=6$, $42.0\%$ at $L=8$ and ${\approx}39.6\%$ at $L=10$. Those contribute nothing
to $w_S$, and they do change $H_S$. The depth uncertainty quoted in Sec.~\ref{sec:frontier} is the
full recomputation of $\mathrm{FR}^{\ast}$ with $L=12$ capped at $n_\ell=300$: it moves by $0.0024$,
$0.0021$, $0.0012$ and $0.0000$ at $\eta=0.10$, $0.15$, $0.18$ and $0.25\,t$.

\paragraph{Reorthogonalization and the certified object.} Full reorthogonalization is used at every
point of the sweep, without which the projector is not a projector: the orthogonality defect is at
most $4.7\times10^{-15}$, and $P_{\mathcal K}\phi=\phi_S$ to $1.1\times10^{-14}$. At $L=12$ the
certificate is evaluated at $n_\ell=300$--$500$ on up to $658\,627$ determinants, the dense $A_S$
never being formed---which is why what the bound controls is $A_{\mathcal K}$ at the declared depth
and not $A_S$ (Sec.~\ref{sec:frontier:hyp}).

\paragraph{Anatomy of the displacement.} The percentages quoted in
Sec.~\ref{sec:frontier:obstruction} are medians over $15$ of the $18$ $(\eta,\text{target})$ cells at
$L=8$; three cells are censored from above, the certificate never reaching the target inside the grid,
and are excluded rather than imputed. Pooling $L=6$ with $L=8$ shifts them to $52\%$, $31\%$, $11\%$
and $4.1\%$, which is why the $L=8$ figures are the ones quoted. One caveat travels with them, and it
is why they are reported as a diagnosis and not as a quantity: of the rungs of the inequality chain
only the last is computable by a user from $(H,S,\phi,\eta)$, so the decomposition is of an object the
user never sees.
